% ============================================================================
% Fichier  : partie2/P02_C04_journaux.tex
% Titre    : Les Journaux
% Auteur   : MINKA MI NGUIDJOÏ Thierry Emmanuel
% ============================================================================
% Sources : RFC 3227 ; NIST SP 800-92 (gestion des journaux) ; loi 2010/012.
% ============================================================================

\chapter{Les Journaux}
\label{chap:journaux}
\chaptermark{Les journaux}

\epigraph{%
    « Un journal que l'on ne collecte pas est une mémoire que l'on efface. »%
}{Principe du cours}

\niveauChapitre{Fondamental}{Chapitre~\ref{chap:pourquoi}~; notions de
systèmes et de réseaux.}

\begin{boxobjectifs}
\begin{itemize}
    \item Identifier les sources de journaux utiles à la détection.
    \item Comprendre la normalisation et l'horodatage.
    \item Fixer une politique de collecte et de conservation.
    \item Relier la journalisation aux obligations légales.
    \item Repérer les angles morts d'une journalisation partielle.
\end{itemize}
\end{boxobjectifs}

\bigskip

% ============================================================================
\section*{Question directrice : parmi tous les journaux de LiZbiZ, lesquels auraient révélé l'intrusion~?}
% ============================================================================

L'attaque a laissé des traces~: une connexion anormale d'un compte de
service, une création de compte, un service installé à distance. Encore
fallait-il que ces journaux existent, soient collectés et conservés. Ce
chapitre recense les sources, leur normalisation et leur conservation.

\medskip
\textbf{Plan du chapitre.}
\begin{enumerate}
    \item Les sources de journaux (section~\ref{sec:c04-sources}).
    \item Normaliser et horodater (section~\ref{sec:c04-normaliser}).
    \item Collecter et conserver (section~\ref{sec:c04-conserver}).
    \item L'obligation de conservation (section~\ref{sec:c04-obligation}).
    \item Les angles morts (section~\ref{sec:c04-angles}).
\end{enumerate}

\separateur

% ============================================================================
\section{Les Sources de Journaux}
\label{sec:c04-sources}
% ============================================================================

\begin{tableaucours}{L{3.6cm}L{8.6cm}}{%
    \tableauHeader{Source} &
    \tableauHeader{Ce qu'elle révèle}
}
Journaux du domaine (Windows) & Connexions, créations de comptes, services, tâches \\
Journaux système (Linux) & Authentifications, sudo, démons, cron \\
Pare-feu et mandataire & Connexions, flux refusés, destinations \\
Serveurs d'application & Accès, erreurs, requêtes \\
Messagerie & Courriels reçus, pièces jointes, signalements \\
Équipements (imprimantes, bornes) & Accès d'administration \\
\end{tableaucours}
\vspace{-0.8em}
\begin{center}{\small\itshape Tableau~4.1. Sources de journaux utiles à la détection}\end{center}

% ============================================================================
\section{Normaliser et Horodater}
\label{sec:c04-normaliser}
% ============================================================================

Des journaux hétérogènes ne se corrèlent que s'ils sont \termecle{normalisés}
(ramenés à un format commun~: qui, quoi, quand, où) et surtout
\termecle{horodatés} de façon cohérente. Une source d'horloge commune (NTP)
est indispensable~: sans elle, impossible de reconstituer une chronologie,
et la preuve perd sa fiabilité.

\begin{boxCRO}
\textbf{Analyse CRO de l'horodatage.} Un horodatage fiable et synchronisé
sert la \cro{0}{+}{+}~: il ne change rien à la confidentialité, améliore la
fiabilité de la trace et son opposabilité devant un tiers. Un horodatage
incohérent ruine l'opposabilité d'une chronologie, même exacte.
\end{boxCRO}

% ============================================================================
\section{Collecter et Conserver}
\label{sec:c04-conserver}
% ============================================================================

\begin{boxprocedure}
\textbf{Politique de collecte des journaux}
\begin{enumerate}
    \item Lister les sources et leur criticité.
    \item Centraliser les journaux hors des machines surveillées (un
          attaquant efface les journaux locaux).
    \item Protéger l'intégrité (accès restreint, écriture seule, condensats).
    \item Synchroniser les horloges (NTP).
    \item Fixer une durée de conservation conforme au droit et aux besoins.
\end{enumerate}
\end{boxprocedure}

% ============================================================================
\section{L'Obligation de Conservation}
\label{sec:c04-obligation}
% ============================================================================

\begin{boxjuris}
La loi n°~2010/012 et la loi n°~2010/013 sur les communications électroniques
imposent aux opérateurs et aux entreprises la conservation des données de
connexion et leur mise à disposition des autorités. Les durées et modalités
précises se lisent dans le texte publié~; le cours en retient le principe~:
des journaux fiables, conservés, sont une obligation autant qu'un outil.
\end{boxjuris}

% ============================================================================
\section{Les Angles Morts}
\label{sec:c04-angles}
% ============================================================================

\begin{boxscenario}{Opération BÉNOUÉ. Ce qui n'était pas journalisé}
\textbf{Constatation.} Les contrôleurs de domaine journalisaient les
connexions, mais le serveur de fichiers \texttt{FS-01} n'envoyait rien au
central, et l'horloge d'un poste de Garoua avançait de 40~minutes.

\medskip
\textbf{Analyse} (confiance élevée). L'intrusion a laissé des traces
visibles sur le domaine, mais l'exfiltration depuis \texttt{FS-01} est restée
invisible faute de collecte, et l'écart d'horloge a brouillé la chronologie.
Deux angles morts à combler en priorité.
\end{boxscenario}

\begin{boxmetier}{Analyste SOC}
\textbf{Cartographier avant de corréler.} L'analyste tient à jour la carte
des sources journalisées et de leurs manques~: une détection n'est fiable que
sur un périmètre connu. Ce qui n'est pas collecté doit être dit, non ignoré.
\end{boxmetier}

\separateur

\begin{boxresume}
\textbf{Ce qu'il faut retenir de ce chapitre~:}
\begin{itemize}
    \item \textbf{Les sources}~: domaine, système, pare-feu, applications,
          messagerie, équipements.
    \item \textbf{Normaliser et horodater} conditionnent la corrélation.
    \item \textbf{Centraliser et protéger} les journaux hors des machines
          surveillées.
    \item \textbf{Conserver} est une obligation légale.
    \item \textbf{Les angles morts} (sources non collectées, horloges
          dérivées) sont des priorités.
\end{itemize}
\end{boxresume}

\begin{boxinfo}
\textbf{Points de contrôle de l'élève.} Je sais~: citer les sources de
journaux~; expliquer normalisation et horodatage~; rédiger une politique de
collecte~; rappeler l'obligation~; repérer un angle mort.
\end{boxinfo}

\bigskip

\begin{boxexo}[title={\ding{45}\quad Exercice~4.1. Sources \niveaubase}]
Quelle source journalise la création d'un compte~? L'exfiltration depuis un
partage de fichiers~?
\end{boxexo}

\begin{boxexo}[title={\ding{45}\quad Exercice~4.2. Horloge \niveaubase}]
Pourquoi une horloge commune (NTP) est-elle indispensable~?
\end{boxexo}

\begin{boxexo}[title={\ding{45}\quad Exercice~4.3. Centraliser \niveaubase}]
Pourquoi centraliser les journaux hors des machines surveillées~?
\end{boxexo}

\begin{boxexo}[title={\ding{45}\quad Exercice~4.4. Normaliser \niveaumoyen}]
Qu'apporte la normalisation à la corrélation de journaux hétérogènes~?
\end{boxexo}

\begin{boxexo}[title={\ding{45}\quad Exercice~4.5. CRO \niveaumoyen}]
Expliquez l'analyse CRO de l'horodatage de l'opération BÉNOUÉ.
\end{boxexo}

\begin{boxexo}[title={\ding{45}\quad Exercice~4.6. Angle mort \niveaumoyen}]
Pourquoi l'exfiltration depuis \texttt{FS-01} est-elle restée invisible~?
\end{boxexo}

\begin{boxexo}[title={\ding{45}\quad Exercice~4.7. Conservation \niveaumoyen}]
Quels critères fixent la durée de conservation des journaux~?
\end{boxexo}

\begin{boxexo}[title={\ding{45}\quad Exercice~4.8. Intégrité \niveaumoyen}]
Comment garantir qu'un attaquant ne puisse pas effacer les journaux
centralisés~?
\end{boxexo}

\begin{boxexo}[title={\ding{45}\quad Exercice~4.9. Politique \niveauexpert}]
Rédigez la politique de journalisation de LiZbiZ~: sources, centralisation,
intégrité, conservation, responsabilités.
\end{boxexo}

\begin{boxexo}[title={\ding{45}\quad Exercice~4.10. Plan de comblement \niveauexpert}]
À partir des angles morts de BÉNOUÉ, proposez un plan de comblement priorisé.
\end{boxexo}

\bigskip

\begin{boxinfo}
\textbf{Pour aller plus loin.}
\begin{itemize}
    \item RFC~3227~\cite{rfc3227}~; loi 2010/012~\cite{loi2010012} et
          2010/013~\cite{loi2010013}.
\end{itemize}
\end{boxinfo}

\begin{boxinfo}
\textbf{Transition.} Collecter ne suffit pas~; il faut corréler~: le
Chapitre~\ref{chap:siem} traite du SIEM.
\end{boxinfo}
